%% =====================================================================
%%  Definition slides -- figures for slides 5, 6, 8, 13 of presentation.md
%%
%%  1. def_polytope.png              -- convex hull of finitely many points
%%  2. def_supporting_example.png    -- f(x,y) = -y, c = 0, H = x-axis
%%  3. def_supporting_hyperplane.png -- general supporting hyperplane
%%  4. def_boundary.png              -- closure / interior / boundary of {|s| <= 1}
%%
%%  Build: pdflatex definitions_figures.tex
%%         pdftoppm -r 300 -png definitions_figures.pdf defs
%% =====================================================================
\documentclass[10pt]{article}
\usepackage{amsmath,amssymb}
\usepackage{tikz}
\usetikzlibrary{arrows.meta,calc,patterns}
\usepackage[active,tightpage]{preview}
\pagestyle{empty}
\setlength{\parindent}{0pt}

\tikzset{
  omega/.style     = {fill=orange!15, draw=orange!70!black, line width=1.5pt},
  face/.style      = {draw=red!78!black, line width=4.5pt, line cap=round},
  supphyp/.style   = {draw=red!78!black, line width=1.8pt, -{Latex[length=4mm,width=3mm]}},
  hyperplane/.style= {draw=red!70!black, line width=1.6pt},
  outside/.style   = {pattern=north east lines, pattern color=red!40, opacity=0.35},
  outsideb/.style  = {pattern=north east lines, pattern color=blue!45, opacity=0.35},
  vertex/.style    = {circle, fill=orange!85!black, draw=white, line width=1pt, inner sep=2.4pt},
  point/.style     = {circle, fill=black, inner sep=2.4pt},
  contact/.style   = {circle, fill=red!80!black, draw=white, line width=1.2pt, inner sep=3pt},
  sub/.style       = {font=\small\sffamily, fill=white, fill opacity=0.95, text opacity=1,
                       rounded corners=2pt, inner sep=3pt},
  subb/.style      = {font=\small\sffamily\bfseries, fill=white, fill opacity=0.95, text opacity=1,
                       rounded corners=2pt, inner sep=3pt},
  tag/.style       = {font=\large\sffamily\bfseries, fill=white, fill opacity=0.92, text opacity=1,
                       rounded corners=3pt, inner sep=4pt},
}

\begin{document}

%% ---------------------------------------------------------------------
%% 1. A polytope is the convex hull of finitely many points.
%%    Every point is a convex combination of the vertices.
%% ---------------------------------------------------------------------
\begin{preview}
\begin{tikzpicture}[x=1cm,y=1cm, line join=round]
  \def\r{2.25}
  % Pentagon vertices v_1 .. v_5 (counter-clockwise).
  \coordinate (v1) at (90:\r);
  \coordinate (v2) at (18:\r);
  \coordinate (v3) at (-54:\r);
  \coordinate (v4) at (-126:\r);
  \coordinate (v5) at (162:\r);
  \coordinate (p)  at (0,0.45);

  \draw[omega] (v1) -- (v2) -- (v3) -- (v4) -- (v5) -- cycle;

  % Each vertex feeds p with its convex weight.
  \foreach \i/\v in {1/v1,2/v2,3/v3,4/v4,5/v5}
    \draw[orange!70!black, line width=0.8pt, dashed, -{Latex[length=2mm,width=1.5mm]}]
      (\v) -- (p);

  \foreach \v in {v1,v2,v3,v4,v5} \node[vertex] at (\v) {};
  \node[point] at (p) {};

  \node[anchor=south west] at (v1) {$v_1$};
  \node[anchor=south west] at (v2) {$v_2$};
  \node[anchor=north]      at (v3) {$v_3$};
  \node[anchor=north east] at (v4) {$v_4$};
  \node[anchor=north east] at (v5) {$v_5$};

  % Labels are kept outside the bundle of arrows.
  \node[sub, anchor=south west] at (0.10,0.52) {$p$};
  \node[tag, text=orange!70!black, anchor=north east] at (-2.30,-1.30) {$\Omega$};
  \node[sub, anchor=north west, align=left] at (2.35,-1.15)
    {$p=\sum_{i=1}^{5}\lambda_i v_i$\\
     $\lambda_i \ge 0,\ \sum_i \lambda_i = 1$};
\end{tikzpicture}
\end{preview}

%% ---------------------------------------------------------------------
%% 2. Example: f(x,y) = -y, c = 0, H = the x-axis.
%%    The pentagon sits above H; the contact set is its bottom edge.
%% ---------------------------------------------------------------------
\begin{preview}
\begin{tikzpicture}[x=1cm,y=1cm, line join=round]
  \coordinate (A) at (-1.6,0);
  \coordinate (B) at ( 2.2,0);
  \coordinate (C) at ( 2.9,1.7);
  \coordinate (D) at ( 0.9,3.1);
  \coordinate (E) at (-1.9,2.0);
  \coordinate (x0) at (0.35,0);

  % The side {f > c} = {y < 0}, below the axis.
  \path[outsideb] (-3.4,0) -- (4.0,0) -- (4.0,-2.5) -- (-3.4,-2.5) -- cycle;

  \draw[omega] (A) -- (B) -- (C) -- (D) -- (E) -- cycle;

  % H = {y = 0} and the exposed set Omega \cap H.
  \draw[hyperplane] (-3.4,0) -- (4.0,0);
  \draw[face] (A) -- (B);

  % f = -y increases downwards (arrow left of pentagon, labels clear of H).
  \draw[supphyp] (-2.90,2.40) -- (-2.90,0.60);
  \node[sub, anchor=south] at (-2.90,2.50) {$f=-\,y$};
  \node[sub, anchor=east] at (-3.05,1.50) {down};

  \node[contact] at (x0) {};

  \node[sub, anchor=west] at (2.50,0.14) {$H:\ y=0$};
  \node[sub, anchor=north] at (0.30,-0.20) {$\Omega\cap H$: exposed edge};
  \node[sub, anchor=west] at (0.58,0.10) {$x_0$};
  \node[tag, text=orange!70!black] at (0.90,2.45) {$\Omega$};
  \node[tag, text=blue!65!black]  at (2.30,-1.85) {$\{f>c\}=\{y<0\}$};
\end{tikzpicture}
\end{preview}

%% ---------------------------------------------------------------------
%% 3. General supporting hyperplane: Omega in the closed half-space {f <= c},
%%    the complementary open side {f > c} shaded, Omega \cap H the flat part
%%    of the boundary in contact with H, x_0 a contact point.
%% ---------------------------------------------------------------------
\begin{preview}
\begin{tikzpicture}[x=1cm,y=1cm, line join=round]
  % The exposed face is the TOP edge A--E, lying on H; every other vertex is
  % strictly below H, so Omega is contained in the closed half-space {f <= c}
  % and the side above H is the complementary open half-space {f > c}.
  \coordinate (A) at (0.6,0.8);      % on H:  y = 0.38235x + 0.57059
  \coordinate (E) at (4.0,2.1);      % on H
  \coordinate (C) at (5.2,1.2);
  \coordinate (D) at (4.4,-0.4);
  \coordinate (F) at (1.6,-0.9);
  \coordinate (G) at (0.2,0.1);
  \coordinate (x0) at (2.3,1.45);

  % The side {f > c}: above the supporting line.
  \path[outside] (-1.4,0.036) -- (6.6,3.095) -- (6.6,6.2) -- (-1.4,6.2) -- cycle;

  \draw[omega] (A) -- (E) -- (C) -- (D) -- (F) -- (G) -- cycle;

  % The supporting hyperplane H and the exposed set Omega \cap H.
  \draw[hyperplane] (-1.4,0.036) -- (6.6,3.095);
  \draw[face] (A) -- (E);

  % grad f, perpendicular to H and pointing into {f > c}.
  \draw[supphyp] (0.0,0.571) -- (-0.54,1.97) node[above, sub] {$f$};

  \node[contact] at (x0) {};
  \node[sub, anchor=west] at (2.42,1.40) {$x_0$: contact point};

  \node[tag, text=orange!70!black] at (2.70,0.45) {$\Omega$};
  \node[tag, text=red!70!black]   at (4.85,5.15) {$\{f>c\}$};
  \node[tag, text=blue!65!black]  at (-0.95,-1.30) {$\{f\le c\}$};
  \node[sub, anchor=west]         at (5.50,1.70) {$H:\ f(x)=c$};
  \node[sub, anchor=south]        at (2.30,1.72) {$\Omega\cap H$: exposed set};
\end{tikzpicture}
\end{preview}

%% ---------------------------------------------------------------------
%% 4. Closure, interior and boundary of S = {s : |s| <= 1}.
%%    A rim point s_0 has a small ball meeting both S and S^c.
%% ---------------------------------------------------------------------
\begin{preview}
\begin{tikzpicture}[x=1.35cm,y=1.35cm, line join=round]
  \def\R{1.7}
  \coordinate (s0) at (60:\R);

  \fill[blue!7] (0,0) circle[radius=\R];
  \draw[blue!70!black, line width=2pt] (0,0) circle[radius=\R];

  % A small ball about s_0 meets S and its complement S^c.
  \draw[red!80!black, dashed, line width=1.1pt] (s0) circle[radius=0.44];
  \node[contact] at (s0) {};
  \node[sub, anchor=south west] at ($(s0)+(0.48,0.32)$) {$s_0\in\partial S$};

  \node[tag, text=blue!70!black]  at (0,0)               {$\operatorname{int}S=\{|s|<1\}$};
  \node[sub, anchor=east]          at (-1.78,1.05)       {$S=\{|s|\le 1\}$};
  \node[sub, anchor=north]         at (0.55,-1.90)      {$S^{c}=\{|s|>1\}$};
  \node[sub, anchor=west]         at ( 1.95,0.95)      {$\partial S=\{|s|=1\}$ (the rim)};
  \node[sub, anchor=west]         at ( 1.95,0.45)      {$B(s_0,\varepsilon)$ meets both};
  \node[sub, anchor=west]         at ( 1.95,0.05)      {$S$ and $S^{c}$};
\end{tikzpicture}
\end{preview}

%% ---------------------------------------------------------------------
%% 5. Exposed edge: Omega = conv{(0,0),(1,0),(1,1),(0,1)} and H = {y = 0},
%%    so Omega \cap H = [(0,0),(1,0)] has affine dimension 1.
%% ---------------------------------------------------------------------
\begin{preview}
\begin{tikzpicture}[x=1.9cm,y=1.9cm, line join=round]
  \coordinate (a) at (0,0);
  \coordinate (b) at (1,0);
  \coordinate (c) at (1,1);
  \coordinate (d) at (0,1);
  \coordinate (m) at (0.5,0);

  % The side {f > c} = {y < 0}, below H.
  \path[outsideb] (-1.75,0) -- (2.85,0) -- (2.85,-1.5) -- (-1.75,-1.5) -- cycle;

  \draw[omega] (a) -- (b) -- (c) -- (d) -- cycle;

  % H = {y = 0} and the exposed edge Omega \cap H.
  \draw[hyperplane] (-1.75,0) -- (2.85,0);
  \draw[face] (a) -- (b);

  \node[contact] at (m) {};
  \node[sub, anchor=south] at (0.5,0.14) {$E$};
  \node[sub, anchor=north] at (0.5,-0.20) {$\Omega\cap H=[\,(0,0),(1,0)\,]$};

  \node[sub, anchor=south east] at (-0.05,0.08) {$H:\ y=0$};
  \node[sub, anchor=west] at (1.15,0.62)
    {$\dim\operatorname{aff}(E)=1$};
  \node[sub, anchor=west] at (1.15,0.32)
    {exposed edge};
  \node[tag, text=orange!70!black, anchor=south] at (0.5,1.28) {$\Omega$};
  \node[tag, text=blue!65!black, anchor=north] at (1.35,-0.55) {$\{f>c\}$};
\end{tikzpicture}
\end{preview}

\end{document}